\section{Results}
\label{sec:results}
We first show that on WikiText103 ALiBi is efficient and enables training models with short input subsequences that outperform strong baselines even when the ALiBi models extrapolate to more than six times the number of tokens that they were trained on. 
We then take the same hyperparameters for our method (the set of slopes) that worked on WikiText-103 and show that -- with no modification -- they provide strong results on a dataset in a very different domain: books. 
Finally, we show that a 1.3B parameter model trained with AliBi on a much larger (461 GB) dataset with much more compute provides a superior alternative to the sinusoidal method since it achieves similar perplexity scores while running faster and using less memory (since it is trained on shorter inputs). 

While multiple alternatives to the position methods presented in~\cite{vaswani} have been proposed, few have been adopted in large (1B or more parameter) LMs since  that setting is much more challenging than the smaller scale experiments. GPT-3 and Jurassic-1~\citep{J1WhitePaper} use the learned position embedding method from Vaswani et al., and GPT-J uses the rotary method.
Our results on the 1.3B parameter model show our method's ability to generalize to larger models, dataset sizes and training durations without retuning the hyperparameter. 


\subsection{Results on WikiText-103 and Toronto BookCorpus}

\begin{figure}[h]
\begin{center}
\includegraphics[width=.65\textwidth]{figures/wt103_extra_lsb_all.pdf} %
\end{center}
\caption{\al models trained and evaluated on varying sequence lengths on the WikiText-103 validation set and the sinusoidal baseline (not evaluated on longer sequences). All of our models outperform the sinusoidal ones even when trained on fewer tokens. Appendix Table~\ref{tab:LSB_wt103} has exact perplexities, more \al models (trained on fewer tokens), and results for rotary and T5 bias models. }
\label{fig:wt103_extra_lsb_all}
\end{figure}


We first develop our method on the WikiText-103 corpus~\citep{pointer}, replacing the sinusoidal position embeddings in the language model of~\cite{baevski} with \al. 


Figure~\ref{fig:wt103_extra_lsb_all} (and the corresponding Appendix Table~\ref{tab:LSB_wt103}) show our results for models trained with varying numbers of input subsequence tokens ($L$), extrapolating to longer subsequence lengths on the validation dataset. 
Our first observation is that, without extrapolation, for every $L$, our models outperform  those using the sinusoidal method, sometimes by a  significant %
amount. For example, the Baevski \& Auli model achieves 18.67$\pm$0.24 (std.~dev.) perplexity when trained with $L = 3072$ input tokens, but our $L=3072$ model achieves 17.60 perplexity (when both models evaluate with \li = 3072).

Our second observation is that all of our models can extrapolate, and they obtain improved perplexity scores when handling more tokens than they observed during training. For example, our model trained on 512 tokens (which achieves 19.73 perplexity when evaluating subsequences of length 512 in the development set) achieves a perplexity score of 18.40 on the development set when extrapolating to subsequences of length 3072. Surprisingly, this surpasses the score that the $L=3072$ sinusoidal model obtains on the development set by a statistically significant margin. 
Note that all our models trained on $L=512$ to $L=2048$ outperform the sinusoidal baseline trained on $L=3072$ when extrapolating to \li = 3072 even though those models all take much less time to train since they train on shorter subsequences (Appendix Figure~\ref{fig:wt103_train_speed_vs_ppl} compares training speed to perplexity for these models)! The $L=512$ model is 1.84 times faster to train and yet still outperforms the $L=3072$ sinusoidal model when extrapolating to $\li = 3072$. In addition, training the $L=3072$ sinusoidal model requires a GPU with more than 16 GB of memory to fit the large attention matrices, which our $L=512$ outperforms even though it can be trained on a GPU with much less memory due to much smaller attention matrices. 

Additionally,  Table~\ref{tab:LSB_wt103} (in the appendix) also shows that, for $L$s  of 1024 and 3072, our method performs better than the rotary and T5 bias models even when \li = $L$ (i.e., no extrapolation is occurring).
Figure~\ref{fig:wt103_extra} (and the corresponding Appendix Tables~\ref{tab:appendix:wt103_extra_512} and~\ref{tab:appendix:wt103_extra_1024}) more broadly explore our method vs.~the other position methods. They show that the T5 bias (the best of the baselines) improves perplexity until \li is around $2L$, but on the WikiText-103 dataset our method continually improves perplexity until at least around $3L$, with the $L=512$ model improving perplexity even when \li exceeds 12k tokens. Even when unable to improve perplexity given longer sequences, \al always maintains strong performance as more tokens are added. 



Appendix Table~\ref{tab:wt103_test} shows that our results on the validation set also transfer to the test set of WikiText-103.
Currently, almost all models that present results on WikiText-103 use sliding window evaluation  (defined in \S\ref{sec:analysis}) to compute perplexities. We apply that method to our (and to the sinusoidal, rotary and T5 bias) models in Appendix Table~\ref{tab:wt103_test_sota}. We find that our L = 3072 model surpasses the performance of Transformer-XL~\citep{transformer-xl}, the Sandwich~\citep{sandwich}, and Shortformer~\citep{shortformer} models. Our results are similar to the ones obtained with staged training~\citep{shortformer} but fall short of results obtained by Routing Transformer~\citep{roy2020efficient} and kNN-LM~\citep{khandelwal20generalization}. The methods used in those models are orthogonal to ours, and we hypothesize that combining them with ours might lead to even larger performance increases. 


After developing our method on WikiText-103, in Appendix Section~\ref{subsec:tbc}, we run one set of experiments on a  different domain (books) using a similar model architecture and without modifying any of the \al hyperparameters (the slopes) and show that our results fully transfer to this new domain. Our models are able to both surpass the sinusoidal baseline when not extrapolating while also outperforming it when extrapolating to longer sequences. 


\subsection{Results on the CC100+RoBERTa Corpus}

Our final set of experiments investigates whether \al transfers to a larger model trained with a larger computational budget on a larger dataset than the ones we previously used. We show that our method achieves strong results in this more challenging setting, obtaining similar performance to the sinusoidal baseline while using significantly less memory, since we train on shorter subsequences. 

The dataset we choose is a combination of the datasets used to train the RoBERTa~\citep{roberta} implementation of BERT~\citep{bert} and the English part of the CC-100 corpus introduced in~\cite{cc-100}, for a total of 461 GB. The RoBERTa training corpus---i.e., the Toronto Book Corpus~\citep{zhu2015aligning}, English Wikipedia, CC-News~\citep{ccnews}, OpenWebText~\citep{openwebtext} and Stories~\citep{stories})---is 161 gigabytes, and the English part of the CC-100 corpus is 300 gigabytes. 
The validation set contains 649K tokens.

 





Our models for this dataset have 25 transformer layers with 16 heads and a dimension of 2048, with an 8192 hidden dimension of the feedforward sublayers.
These models have 1.3B parameters. 
We train our models for one epoch, which is 50k updates on 128 V100 GPUs.



\begin{figure}
\centering
\begin{subfigure}{.45\textwidth} %
  \centering
  \includegraphics[width=\linewidth]{figures/ccr-train-512.pdf}
 
\end{subfigure}%
\begin{subfigure}{.45\textwidth} %
  \centering
 \includegraphics[width=\linewidth]{figures/ccr-train-1k.pdf}
\end{subfigure}
\caption{ On the left (right), a 1.3B-parameter ALiBi model trained on 512 (1024) and evaluated on 1024 (2048) tokens during training, compared to the sinusoidal baseline trained on 1024 (2048) tokens. The \al models obtain strong results even though they use 6\%-11\% less memory since they train on shorter sequences. Appendix Table \ref{tab:ccr} shows memory use and end-of-training perplexities.}
\label{fig:ccr_train_all}
\end{figure}


In Figure~\ref{fig:ccr_train_all} (left), we compare the validation perplexity for \li = 1024 throughout the training process for an \al model trained with \lt = 512 compared to the sinusoidal model trained with \lt = 1024. Since our model is trained on shorter sequences, it is 7\% faster and uses 1.6 GB less memory. We halt training of the sinusoidal baseline when our model reaches the end of its training (one epoch). 
At that time, our model is just 0.06 perplexity away from the baseline even though it was trained on sequences that are half the length of those the baseline used and requires less memory.

In Figure~\ref{fig:ccr_train_all} (right), results become even more impressive, showing that our model trained on \lt = 1024 outperforms by 0.09 perplexity the sinusoidal model trained on \lt = 2048 (when evaluating with \li = 2048) even though our model uses 3.1 GB less memory. Our model maintains a lead in perplexity over the sinusoidal model during the entire training process. By sampling five evenly distributed points across the training process, we compute that our \lt = 1024 model reaches a given perplexity value, on average, 11\% faster than the sinusoidal model does. 

Since our models in these comparisons use much less memory, they allow for stacking more layers, which would further improve performance (with negligible, if any, runtime cost). To keep our experiments as straightforward as possible, however, we do not add layers to our models.

Appendix Table~\ref{tab:ccr_all_50k} presents additional results comparing our models to the sinusoidal baseline when both are trained on the same $L$, showing that \al performs similarly to the sinusoidal baseline when not extrapolating. This contrasts with the  results presented on the smaller datasets, where \al consistently outperforms other position methods even when not extrapolating, suggesting that ALiBi's inductive bias provides additional benefits for lower-resource language modeling.


\begin{figure}[h]
\centering
\begin{subfigure}{.45\textwidth}  %
  \centering
  \includegraphics[width=\linewidth]{figures/ccr_extra_512.pdf}
\end{subfigure}%
\begin{subfigure}{.45\textwidth} %
  \centering
  \includegraphics[width=\linewidth]{figures/ccr_extra_1024.pdf}
\end{subfigure}
\caption{The ALiBi and sinusoidal models (with both $L$ = 512 and 1024) trained for 50k updates (1 epoch) on the CC100+RoBERTa corpus, extrapolating on the validation set. %
\al achieves the best results at around $2L$ but maintains strong performance even up to 10000 tokens in these experiments.}
\label{fig:ccr_extra}
\end{figure}


Figure~\ref{fig:ccr_extra} shows that our models trained on $L=512$ and $L=1024$ achieve the best results when extrapolating to about double the tokens that they were trained on.
Specifically, the \lt = 512 model (that obtains 9.79 perplexity when \li = 512) achieves its best score (9.3) when extrapolating to 1012 tokens, and the \lt = 1024 model (that obtains 9.16 perplexity when \li = 1024) achieves its best score (8.9) when extrapolating to 2024 tokens. 


One possible explanation is that the subsequences the model observes during training are up to $L$ tokens long. When performing inference on subsequences of length $2L$, half of the subsequences the model consumes are as long as the examples seen during training. When inference is performed on subsequences of length $2L+1$ or longer, less than half of the predictions the model makes are on subsequences of lengths seen during training, and that might degrade performance. 

The sinusoidal model cannot extrapolate at all in this setting, with its performance degrading for both the \lt = 512 and 1024 models as soon as one token more than $L$ is added during evaluation. 


In Appendix \ref{sec:analysis}, we find that \al's  edge over sinusoidal embeddings is largely explained by its improved avoidance of the early token curse.  We posit  that future work building on \al might achieve further gains by more efficiently  exploiting longer histories.

